\documentclass[aspectratio=169]{beamer}

% ---------- ENCODING & FONTS ----------
\usepackage[utf8]{inputenc}
\usepackage[T1]{fontenc}

% ---------- THEME & LOOK ----------
\usetheme{Madrid}
\usecolortheme{seahorse}
\usefonttheme{professionalfonts}
\setbeamertemplate{navigation symbols}{}
\setbeamercovered{transparent}
\setbeamertemplate{itemize items}[circle]
\setbeamertemplate{enumerate items}[default]
\setlength{\parskip}{0.25em}

% ---------- COLORS ----------
\usepackage{xcolor}
\definecolor{accent}{RGB}{0,92,184}
\definecolor{soft}{RGB}{233,244,255}
\setbeamercolor{structure}{fg=accent}
\setbeamercolor{title}{fg=white,bg=accent}
\setbeamercolor{frametitle}{fg=accent}
\setbeamercolor{block title}{bg=accent!12,fg=accent}
\setbeamercolor{block body}{bg=soft}

% ---------- PACKAGES ----------
\usepackage{booktabs}
\usepackage{amsmath, amssymb}
\usepackage{array}
\usepackage{multirow}
\usepackage{tabularx}
\usepackage{hyperref}
\hypersetup{colorlinks=true,linkcolor=accent,urlcolor=accent}

% ---------- SAFE TRANSITION FALLBACK ----------
\makeatletter
\@ifundefined{transfade}{\newcommand{\transfade}[1][]{}}{}
\makeatother

% ---------- PLACEHOLDER BOX (compile-safe; no images required) ----------
\newcommand{\placeholderbox}[2][0.82\linewidth]{%
  \begin{center}
    \fbox{\begin{minipage}[c][0.46\textheight][c]{#1}
      \centering \vspace{0.25em}
      \textbf{#2}\\[0.6em]
      \emph{Insert figure/diagram here later.}
    \end{minipage}}
  \end{center}
}

% ---------- TITLE ----------
\title{\textbf{Union Impact on Wage and Nonwage Outcomes  \vspace{.5cm}\\  }}
\subtitle{ {\large ECON 300 - Labour Economics}}
%\institute{UNBC}
\author{Muhebullah Karimzada}

\date{Winter 2026}



\begin{document}

% ============================================================
% TITLE
% ============================================================
\begin{frame}[plain]
  \titlepage
\end{frame}

% ============================================================
\section{Learning Objectives \& Main Questions}
% ============================================================

\begin{frame}{Learning Objectives}
\transfade
\begin{itemize}[<+->]
  \item Discuss difficulties in measuring the impact of unions on members' wages and methods to overcome them.
  \item Describe mechanisms whereby unions may affect nonunion wages.
  \item Explain how unions affect wage inequality and the distribution of income.
  \item Evaluate how union impacts differ by skill, gender, sector, and product-market competition.
  \item Analyze effects on productivity, investment, profitability, and overall efficiency.
\end{itemize}
\end{frame}

\begin{frame}{Main Questions}
\transfade
\begin{itemize}[<+->]
  \item How do we estimate the \alert{causal} wage effect of unions?
  \item By how much do unions raise members' wages, and for whom the most?
  \item What happens to nonunion wages and employment?
  \item Do unions raise or reduce overall income inequality?
  \item What are the effects on productivity and firm profits?
\end{itemize}
\end{frame}

% ============================================================
\section{Union Wage Impact: Two-Sector Model}
% ============================================================

\begin{frame}{Basic Two-Sector Framework}
\transfade
\begin{itemize}[<+->]
  \item \textbf{Union--nonunion wage differential:} percentage difference in wages between union and nonunion workers.
  \item Predictions:
    \begin{itemize}[<+->]
      \item Union increases wage in the organized sector, reduces employment there.
      \item Displaced workers flow to nonunion sector $\Rightarrow$ wage there falls.
    \end{itemize}
\end{itemize}
\end{frame}

\begin{frame}{Figure 15.1 -- Two-Sector Model of General Equilibrium \textit{(Placeholder)}}
\transfade
\placeholderbox{Figure 15.1 -- Two-Sector Model of General Equilibrium.}
\end{frame}

\begin{frame}{Additional Impacts and Threat Effects}
\transfade
\begin{itemize}[<+->]
  \item Nonunion firms may raise wages to compete for workers.
  \item \textbf{Threat effect:} nonunion employers raise wages to reduce risk of unionization.
  \item Threat effect stronger when:
    \begin{itemize}[<+->]
      \item Union wage is high
      \item Nonunion sector easy to organize
      \item Potential union aggressive
      \item Employers strongly averse to unionization
    \end{itemize}
\end{itemize}
\end{frame}

\begin{frame}{Wait or Queue Unemployment}
\transfade
\begin{itemize}[<+->]
  \item Workers may rationally \textbf{wait} for union jobs due to higher pay.
  \item Nonunion supply increase $<$ union employment reduction (queuing attenuates spillovers).
\end{itemize}
\end{frame}

% ============================================================
\section{Measurement \& Empirical Evidence}
% ============================================================

\begin{frame}{Estimating the Pure Union Differential}
\transfade
\begin{itemize}[<+->]
  \item Must control for:
    \begin{itemize}[<+->]
      \item Skills and labour quality
      \item Job assignment characteristics
      \item Education, training, experience
      \item Other productivity-related factors
    \end{itemize}
\end{itemize}
\end{frame}

\begin{frame}{Empirical Evidence: Headline Findings}
\transfade
\begin{itemize}[<+->]
  \item Unions cause higher wages.
  \item Typical union--nonunion differential:
    \begin{itemize}[<+->]
      \item U.S.: 10--20\%
      \item Canada: 10--25\%
    \end{itemize}
  \item \textbf{Longitudinal} (panel) estimates controlling for person fixed effects: $\sim$10\%.
\end{itemize}
\end{frame}

\begin{frame}{Worked Example -- Computing a Union Premium}
\transfade
\small
\begin{itemize}[<+->]
  \item Nonunion pay structure (illustrative): $Y = 8000 + 2000\,S + 500\,X$.
  \item For a union worker with $S=12$ and $X=10$:
\[
Y_U^\ast = 8000 + 2000(12) + 500(10) = \$37{,}000.
\]
  \item Observed union earnings: $Y_U = \$44{,}000$.
  \item \textbf{Pure union premium:} $Y_U - Y_U^\ast = \$7{,}000$.
\end{itemize}
\end{frame}

\begin{frame}{Table: Worked Example Summary (Exact Values)}
\transfade
\small
\begin{tabular}{lcc}
\toprule
\textbf{Quantity} & \textbf{Expression} & \textbf{Value} \\
\midrule
Expected nonunion earnings for union worker & $Y_U^\ast = 8000 + 2000(12) + 500(10)$ & \$37{,}000 \\
Observed union earnings & $Y_U$ & \$44{,}000 \\
Pure union earnings premium & $Y_U - Y_U^\ast$ & \$7{,}000 \\
\bottomrule
\end{tabular}
\end{frame}

\begin{frame}{Decomposition Idea}
\transfade
\begin{itemize}[<+->]
  \item Use Blinder--Oaxaca to split average union--nonunion earnings difference into:
    \begin{itemize}[<+->]
      \item Part due to \textbf{endowments} (education, experience, occupation mix, etc.)
      \item Part due to \textbf{coefficients} (structure or returns differences)
    \end{itemize}
\end{itemize}
\end{frame}

% ============================================================
\section{Heterogeneity in Impacts}
% ============================================================

\begin{frame}{Variation in the Union Wage Impact}
\transfade
\begin{itemize}[<+->]
  \item \textbf{By union density:} larger differentials where organization is widespread; often nonlinear.
  \item \textbf{By firm size:} positive relation with wages in both union and nonunion sectors; stronger for nonunion.
  \item \textbf{By industry and sector:} smaller in public sector; lower where foreign competition is stronger.
  \item \textbf{By skill and gender:} highest among low-skilled; union and nonunion wages similarly responsive to gender.
  \item \textbf{Over cycle/time:} union wages exhibit less cyclical variation.
\end{itemize}
\end{frame}

\begin{frame}{Figure 15.2 -- Union vs.\ Nonunion Wages by Characteristics \textit{(Placeholder)}}
\transfade
\placeholderbox{Figure 15.2 -- Relationship between union and nonunion wages and personal/job characteristics.}
\end{frame}

% ============================================================
\section{Inequality, Dispersion, and Distribution}
% ============================================================

\begin{frame}{Unions and Wage Dispersion}
\transfade
\begin{itemize}[<+->]
  \item Less wage dispersion among union workers in Canada.
  \item Larger impact on blue-collar wages.
  \item Net effect tends to \textbf{reduce} overall wage dispersion and inequality.
\end{itemize}
\end{frame}

\begin{frame}{Figure 15.3 -- Wage Dispersion in Union vs.\ Nonunion Sectors \textit{(Placeholder)}}
\transfade
\placeholderbox{Figure 15.3 -- Wage dispersion comparison between union and nonunion sectors.}
\end{frame}

% ============================================================
\section{Resource Allocation \& Efficiency}
% ============================================================

\begin{frame}{Union Impact on Resource Allocation and Welfare}
\transfade
\begin{itemize}[<+->]
  \item Wage and employment changes affect allocation of labour and capital.
  \item Potential \textbf{deadweight losses} if workers shift to lower-productivity nonunion jobs.
  \item Magnitude depends on:
    \begin{itemize}[<+->]
      \item Labour supply elasticity
      \item Extent of queuing/unemployment
      \item Whether wage--employment are jointly negotiated (efficient contracts)
    \end{itemize}
\end{itemize}
\end{frame}

\begin{frame}{Figure 15.4 -- Union Wage Impact and Allocative Inefficiency \textit{(Placeholder)}}
\transfade
\placeholderbox{Figure 15.4 -- Areas of deadweight loss under two-sector model with union wage.}
\end{frame}

% ============================================================
\section{Nonwage Outcomes}
% ============================================================

\begin{frame}{Fringe Benefits and Total Compensation}
\transfade
\begin{itemize}[<+->]
  \item Nonwage benefits are a higher share of compensation in union firms.
  \item Worker preferences (median voter) and life-cycle demands shape packages.
  \item Benefits can raise retention and reduce turnover.
\end{itemize}
\end{frame}

\begin{frame}{Employment, Work Rules, and Conditions}
\transfade
\begin{itemize}[<+->]
  \item Employers may substitute away from higher-cost union labour; output prices may rise.
  \item Unions may bargain on employment levels; \textbf{featherbedding} can fix labour--capital ratios.
  \item Nonunion settings: more flexibility and fewer rules; faster pace often noted.
\end{itemize}
\end{frame}

\begin{frame}{Turnover, Productivity, Profitability, and Investment}
\transfade
\begin{itemize}[<+->]
  \item Turnover is typically lower in unionized firms (deferred compensation more prevalent).
  \item Unions can increase capital intensity and select more productive labour.
  \item Potential positive channels: morale, communication, and cooperation.
  \item On balance, profit impacts depend on product-market conditions and bargaining outcomes.
\end{itemize}
\end{frame}

% ============================================================
\section{Chapter Summary}
% ============================================================

\begin{frame}{Summary}
\transfade
\begin{itemize}[<+->]
  \item Two-sector model clarifies wage, employment, and spillover effects.
  \item Identification requires rich controls; panel methods shrink differentials toward $\sim$10\%.
  \item Union impacts vary by density, firm size, industry, sector, skill, and competition.
  \item Unions generally compress wage dispersion; efficiency losses hinge on elasticities and contracting.
  \item Nonwage outcomes include higher benefits, lower turnover, and changes in productivity and profits.
\end{itemize}
\end{frame}

% ============================================================
\begin{frame}
\centering
\Huge{\textbf{Thank You!}} \\

\end{frame}

\end{document}
